\begin{abstract}
We characterize the optimal accuracy of pure differentially private pointwise treatment-effect estimation and honest interval inference from independent binary observational records. The model has a scalar covariate on \([0,1]\), an unknown measurable density between \(1/2\) and \(3/2\), treatment probabilities between \(1/4\) and \(3/4\), propensity and control regressions of H\"older order \(1/10\), and a Lipschitz treatment contrast, all with regularity radius \(3\). Within this model, consistency and conditional exchangeability identify the causal contrast almost surely; continuity and positive design density identify its value at the fixed covariate location \(1/2\). For sample size \(n\ge2\) and pure replacement privacy budget \(0<\epsilon\le1\), both minimax expected absolute error and minimax maximal expected length of connected intervals with uniform finite-sample \(90\%\) coverage have order
\[
\min\left\{1,\max\left\{
n^{-1/4},(n^2\epsilon)^{-1/7},(n\epsilon)^{-1/2}
\right\}\right\}.
\]
A count-weighted within-cell pair ratio attains these orders through one joint Laplace release and public tuning. Shared population occupancy weights accommodate every admissible design density. Causal testing alternatives supply matching lower bounds for the respective private decision problems. Along admissible privacy-budget sequences, vanishing optimal risk and honest expected length are equivalent to \(n\epsilon\to\infty\).
\end{abstract}

\section{Introduction}

This paper establishes matching finite-sample orders for private estimation and honest interval inference for a treatment effect at a fixed covariate value. The experiment consists of independent binary observational records with a scalar covariate, fixed treatment overlap, a bounded positive design density, rough nuisance regressions, and a Lipschitz treatment-effect contrast. Within this experiment, a single joint private release attains both the minimax expected absolute-error order and the optimal order of maximal expected length of connected intervals with uniform \(90\%\) coverage. The matching conclusions concern the dependence on sample size and privacy budget, with explicit numerical constants.

Pointwise treatment effects describe how the expected response to treatment varies across observed characteristics. Under consistency and conditional exchangeability, observational treatment-arm regressions identify the conditional potential-outcome contrast. Pointwise evaluation additionally requires a choice of regression version. In the class studied here, contrast continuity and positive covariate density determine a unique continuous representative, making its value at the public evaluation location an identified feature of the observed-record law, as established in \cref{obj:lem:point-version}. Privacy then constrains how that feature can be estimated and reported from confidential records.

The statistical class in \cref{obj:def:complete-model} fixes the covariate domain to \([0,1]\), the evaluation location to \(1/2\), and a common regularity radius of \(3\). It combines a measurable density bounded above and away from zero, treatment overlap, one-tenth H\"older propensity and control regressions, and a Lipschitz contrast for binary outcomes. Releases satisfy pure central differential privacy under replacement adjacency on the full ordered-dataset space.

The smooth contrast permits the two treatment-arm means to share rough local variation that cancels in their difference, while rough treatment selection remains an additional nuisance. This scalar binary experiment makes the interaction between nuisance roughness, local occupancy, and privacy explicit.

Let \(n\) denote sample size and \(\epsilon\) the privacy budget. For every integer \(n\ge2\) and \(0<\epsilon\le1\), define the accuracy benchmark
\[
r(n,\epsilon)=
\min\left\{1,\max\left\{
n^{-1/4},(n^2\epsilon)^{-1/7},(n\epsilon)^{-1/2}
\right\}\right\}.
\]
\Cref{obj:thm:matched-risk-frontier} establishes that minimax maximal expected absolute error is bounded above and below by numerical multiples of this benchmark. \Cref{obj:thm:sharp-interval-frontier} establishes the same order for minimax maximal expected length among private connected-interval procedures satisfying uniform finite-sample \(90\%\) coverage. Both criteria integrate sampling variation and release randomization. The lower bounds range over all everywhere-defined private randomized procedures admitted by the respective criteria.

The attaining construction uses a public window around the evaluation location, partitioned into equal cells. Within each cell, treatment--outcome pair differences form a numerator and treatment pair differences form a denominator. Count weighting gives their population means the same occupancy weights, so the population ratio preserves the cellwise approximation to the target across every admissible measurable design density. The joint statistic has uniformly bounded replacement sensitivity, including changes involving cells with fewer than two records. Independent Laplace perturbations, a public denominator floor, and clipping therefore yield an everywhere-defined private ratio. Public tuning balances localization bias, within-cell nuisance approximation, pair variation, and privacy noise. The uniform squared-error envelope in \cref{obj:thm:uniform-private-upper} supplies an honest interval radius, and candidate-value inversion in \cref{obj:def:interval-handle} reproduces that interval from the same joint release.

The converse uses causal alternatives satisfying the same model restrictions. Shared-sign alternatives vary the nuisance functions while maintaining a common treatment-effect separation; their mixture permits likelihood and component-coupling comparisons. Direct localized alternatives provide the complementary privacy comparisons. The finite testing priors in \cref{obj:lem:frontier-testing-priors} combine these experiments across budget regimes. For scalar estimation, accurate recovery of separated targets would distinguish their released-output laws. For honest intervals, transferred coverage probabilities force simultaneous containment of separated targets with positive probability, and connectedness converts that containment into a length lower bound. This interval argument establishes the expected-length converse directly.

The closest pointwise comparison is \citet[][Theorems 1--2 and Remarks 2 and 11]{KennedyBalakrishnanRobinsWasserman2024}. At dimension one, contrast smoothness one, and average nuisance smoothness one tenth, their rate specializes to \(n^{-1/4}\), with the nuisance average below the relevant one-sixth threshold. For the binary bounded-positive-density experiment here, shared occupancy weighting attains this statistical scale together with the privacy scales \((n^2\epsilon)^{-1/7}\) and \((n\epsilon)^{-1/2}\). The reusable contribution is the combination of density-robust occupancy-weighted attainment with matching private testing comparisons; the connected-interval converse then transfers that frontier directly to honest expected length. \Cref{sec:related-work} gives the detailed comparison of experiments, assumptions, and adjacent private and inferential literature.

The benchmark gives an order-based guide to sample-size and budget choices. As the budget decreases, its dominant term moves from statistical variation to sparse-pair privacy noise, then dense-occupancy privacy noise, and finally bounded-range saturation. \Cref{obj:prop:regime-and-sanity} identifies the transition budgets and establishes that vanishing private minimax absolute risk along admissible budget sequences is equivalent to \(n\epsilon\to\infty\). The interval comparison gives the same condition for vanishing optimal honest expected length.



The paper proceeds through related work, the observation model and decision criteria, and the within-cell private release. It then presents the matching estimation and honest interval frontiers, followed by discussion and limitations. The appendices develop identification and upper-bound ingredients, lower-bound experiments and component coupling, testing reductions, and proofs of the main results.
